\section{Introduction}
\label{sec:introduction}

Vision-language-action models (VLAs) fine-tune large language models on robot demonstrations~\cite{brohan2023_rt2,kim2024_openvla}, under the assumption that fine-tuning injects action-relevant representations.
Yet this assumption has not been tested at the representational level: where do the representations driving action prediction originate, and how does the LLM transform them?
The answers determine whether VLA design should prioritize scaling the LLM or improving vision encoder quality.

Recent work suggests VLAs rely disproportionately on visual input~\cite{grant2026}, with language instructions largely ignored by the action pathway~\cite{haon2025}.
However, these behavioral and single-component analyses have not provided a systematic \emph{representational decomposition} quantifying how each architectural stage independently contributes to action-relevant structure, leaving open whether fine-tuning creates new action knowledge or reconfigures inherited geometry.

We address this gap by decomposing OpenVLA-7B~\cite{kim2024_openvla,karamcheti2024}---a VLA chaining a frozen DinoSigLIP vision backbone~\cite{zhai2023_siglip,oquab2024_dinov2}, an MLP projector, and Llama-2 7B---extracting representations from the projector, each of 32 LLM layers (at both image-token and readout-token positions), and two external baselines (CLIP~\cite{radford2021_clip} and DINOv2~\cite{oquab2024_dinov2}).
We cross training status (trained vs.\ untrained, reinitializing 6.74B LLM parameters) with instruction presence in a $2\times 2$ factorial design, and quantify action-relevant structure using RSA~\cite{kriegeskorte2008} for geometric alignment and MLP probes~\cite{alain2017} for functional decodability.

Three findings emerge from this decomposition.
\emph{First}, the frozen vision encoder provides the dominant source of action-relevant geometry, with self-supervised pretraining exceeding vision-language pretraining.
Standalone DINOv2 achieves the highest action-space alignment (RSA $\rho = 0.692$), far exceeding CLIP ($\rho = 0.557$) and the VLA's own peak internal representation ($\rho = 0.628$ at L12).
Functional probes corroborate this ranking: DINOv2 ($R^2 = 0.634$) $\geq$ Projector ($0.630$) $>$ CLIP ($0.587$).
\emph{Second}, the LLM applies a non-monotonic transformation: the first 12 layers \emph{enhance} geometric alignment above the vision encoder input ($\rho$: $0.612 \to 0.628$), while the remaining 20 layers degrade it ($0.628 \to 0.422$), and the readout token concentrates a functionally superior representation (probe $R^2 = 0.678$ vs.\ $0.630$ for the projector).
\emph{Third}, training produces at most a modest geometric gain (RSA $\Delta = 0.017$--$0.023$ across two benchmarks, $3$--$4\%$ of baseline) while substantially improving functional capacity ($R^2$ $\Delta = +0.015$ to $+0.069$, $p \leq 0.006$)---redistribution dominates over creation.

Our contributions are: (i)~the first per-layer representational decomposition of a VLA using both geometric and functional metrics; (ii)~trained-vs-untrained factorial controls isolating the LLM's contribution from pretrained vision geometry; and (iii)~multi-metric evidence that action representations are inherited from vision encoders and non-monotonically transformed by the LLM, suggesting that vision encoder quality is the primary design lever (\S\ref{sec:discussion}).
